\chapter{Repérage, distances et lieux géométriques}
\section*{Exercice 7.1 -- Distance et milieu}
Dans un repère orthonormé, $A(-1,2)$ et $B(5,10)$. Calculer $AB$ et les coordonnées du milieu $I$.

\section*{Exercice 7.2 -- Triangle rectangle}
Pour $A(0,0)$, $B(4,0)$ et $C(0,3)$, calculer les longueurs et justifier que le triangle est rectangle.

\section*{Exercice 7.3 -- Cercle}
Le cercle a pour équation $(x-2)^2+(y+1)^2=25$. Donner centre et rayon, puis déterminer si $(5,3)$ est sur ce cercle.

\section*{Exercice 7.4 -- Médiatrice}
Déterminer le lieu des points équidistants de $A(-2,0)$ et $B(4,0)$.

\section*{Exercice 7.5 -- Distance à un axe}
Pour $M(3,-4)$, donner les distances aux deux axes. Justifier que la distance à l'axe horizontal est minimale au point $(3,0)$.

\section*{Exercice 7.6 -- Rectangle général}
Prouver par coordonnées que les diagonales d'un rectangle ont même milieu et même longueur. Garder ses côtés $a,b>0$ comme paramètres.
